\section{Part 2: the paper's BERT and Square shapes (\iss{71})}
The paper (arXiv 2603.17800v1, the user's published design) picks its tile from a 1-core microkernel heatmap, whose best on the BPI was 16x15, and runs its GEMM on one core: S1--S5 at 32x7, BERT at 32x6, 16x15 and 32x7. \iss{71} asks which tile wins per shape at 1 and 8 threads in this repo's current design. The first table lists the paper's own GFLOPS (Fig.~11, its own xDSL code and protocol) beside ours for reference; they are not a matched comparison.

\begin{figure}[htbp]\centering
\begin{tikzpicture}
\begin{axis}[house, xbar, bar width=6pt, width=13cm, height=15.5cm, y dir=reverse,
  ytick=data, yticklabels from table={\G/paper_tiles.csv}{label}, enlarge y limits={abs=0.6},
  xmin=0, xmax=140, xtick={0,10,...,90}, xlabel={GOPs by mean time, 8 threads},
  nodes near coords, point meta=explicit symbolic, every node near coord/.append style={font=\scriptsize, anchor=west},
  legend style={at={(0.5,1.02)}, anchor=south, legend columns=-1}]
\addplot[fill=his, draw=none] table[y=idx, x=a8_168, meta=l168]{\G/paper_tiles.csv};
\addplot[fill=vec, draw=none] table[y=idx, x=a8_327, meta=l327]{\G/paper_tiles.csv};
\legend{16x8 (8T winner), 32x7 (paper)}
\end{axis}
\end{tikzpicture}
\note{\iss{71}: jbpi, 8T cores 0--7, uninstrumented, reference protocol, \texttt{--tmin 2}, \texttt{--no-nc-overrides}, vector packer, 19 tiles (the YOLO list, not the paper's full heatmap), 3 launches interleaved with the 1T launches (1T 8T 8T 1T 1T 8T). Core 0 takes the board's interrupts. Bar label: tile, then the cache blocking that run used (from its log).}
\end{figure}

\begin{center}\small
\begin{tabular}{@{}l l l l@{\hspace{1.5em}}l l@{}}
\toprule
 & & \multicolumn{2}{c}{this design, jbpi} & \multicolumn{2}{c}{paper, Fig.~11 (BPI-F3, 1 core)} \\
\cmidrule(lr){3-4}\cmidrule(l){5-6}
 & & \multicolumn{1}{c}{1T (core 4)} & \multicolumn{1}{c}{8T (cores 0--7)} & GEMM+xDSL & OpenBLAS \\
shape & m $\times$ n $\times$ k & best tile, GOPs & best tile, GOPs & tile, GFLOPS & GFLOPS \\
\midrule
B1 & 1024$\times$384$\times$1024 & 32x7 8.06 & 16x8 47.44 & 32x6 7.4 & 6.0 \\
B2 & 384$\times$384$\times$64 & 16x15 12.58 & 16x12 81.99 & 16x15 12.2 & 5.1 \\
B3 & 64$\times$384$\times$384 & 32x7 7.01 & 8x24 27.78 & 16x15 8.5 & 7.3 \\
B4 & 4096$\times$384$\times$1024 & 32x7 8.12 & 16x8 49.33 & 32x7 7.5 & 5.9 \\
B5 & 1024$\times$384$\times$4096 & 32x7 7.80 & 16x8 43.64 & 32x6 6.8 & 5.9 \\
S1 & 1000$\times$1000$\times$1000 & 32x7 8.84 & 16x8 44.70 & 32x7 8.3 & 6.1 \\
S2 & 2000$\times$2000$\times$2000 & 32x7 9.28 & 16x8 47.33 & 32x7 8.5 & 6.8 \\
S3 & 3000$\times$3000$\times$3000 & 32x7 9.54 & 16x8 49.76 & 32x7 8.5 & 6.8 \\
S4 & 4000$\times$4000$\times$4000 & 32x7 9.40 & 16x8 50.99 & 32x7 8.3 & 6.9 \\
S5 & 5000$\times$5000$\times$5000 & 32x7 9.66 & 16x8 46.83 & 32x7 8.4 & 7.0 \\
\bottomrule
\end{tabular}
\note{GOPs, average of 3 launches (\iss{71} \texttt{report/tile\_search\_\{1t,8t\}\_3launches.md}). jbpi, 1T core 4, 8T cores 0--7, uninstrumented, reference protocol, \texttt{--tmin 2}, vector packer. Paper columns: arXiv 2603.17800v1 Fig.~11 (bottom), BPI-F3 (same K1 chip), 1 core, the paper's own xDSL code, toolchain and protocol, tile as labelled in the figure; the figure prints no values, so they are read from the bar heights (pixel-measured, $\pm$0.1). Shown beside ours, not as a matched comparison.}
\end{center}

\begin{center}\small
\begin{tabular}{@{}l rr@{\hspace{2em}}l rr@{}}
\toprule
\multicolumn{3}{@{}c}{1T, one pass over the 10 shapes} & \multicolumn{3}{c@{}}{8T, one pass over the 10 shapes} \\
tile & s & GOPs & tile & s & GOPs \\
\midrule
32x7 & 48.01 & 9.53 & 16x8 & 9.47 & 48.27 \\
32x6 & 51.78 & 8.83 & 16x15 & 10.17 & 44.96 \\
24x9 & 53.28 & 8.58 & 16x12 & 10.22 & 44.74 \\
16x6 & 53.90 & 8.48 & 32x7 & 10.38 & 44.07 \\
24x8 & 54.83 & 8.34 & 16x14 & 10.47 & 43.67 \\
\bottomrule
\end{tabular}
\note{\iss{71} single-tile rankings, same runs and settings. 16x15 over the 1T pass: 7.47 GOPs.}
\end{center}

\begin{center}\small
\begin{tabular}{@{}l l lllll@{}}
\toprule
shape & m $\times$ n $\times$ k & 32x7 & 16x15 & 16x8 & 16x12 & 8x24 \\
\midrule
\csname @@input\endcsname \D/paper_blk_rows.tex
\bottomrule
\end{tabular}
\note{Cache blocking mc / nc / kc of every tile named in this part, per shape, as printed in the run logs (\iss{71}, \iss{83}). It depends on shape and tile only: the same at 1T and 8T and with either packer. 32x7 and 16x15 share their blocking on every shape; nc $<$ n from S2 up.}
\end{center}

\begin{itemize}\setlength\itemsep{1pt}
\item \textbf{1 thread: the paper's picks hold among our 19 tiles.} 32x7 is best on 9 of 10 shapes, 16x15 on B2 (12.58 against 11.76 for 32x7).
\item \textbf{8 threads: 16x8 wins 8 of 10 shapes} and the pass (48.27 GOPs; the paper's 32x7 is 4th at 44.07). B2: 16x12 (81.99; 16x15 81.12). B3: 8x24.
\item \textbf{The paper's 32x7 at 8T} is 8--13\% slower than the best on B1, B4, B5, 28\% on B2, 40\% on B3, and 1--15\% on S1--S5.
\item \textbf{Slow reps at 8T:} S1 and S2 (44.70 and 47.33 GOPs) hold a few much slower reps per launch that pull the mean time down (the cause, likely core 0, is inferred from the YOLO runs, \iss{79} B6).
\item 8T over 1T on the best single tile: 48.27 / 9.53 = 5.1$\times$. Why the winner changes with threads is not measured.
\end{itemize}
\noindent\textit{Limits:} the 19 tiles are the YOLO candidate list (\iss{59}), not the paper's full heatmap; S5 (5000$\times$5000$\times$5000) gets few reps per launch at 1T with \texttt{--tmin 2}.
